% Point-by-point response to the editorial decision on
% "Structure-Aware Row Scaling for Block-Structured Mixed-Integer Programs"
% (submitted as "Generator-Aware Positive Row Scaling for Block-Structured
% Engineering Mixed-Integer Programs"), IJNME.
%
% Build: pdflatex response-letter
% NOTE: the \reviewerquote{} blocks paraphrase each comment; before submission,
% replace them with the verbatim text from the decision letter.

\documentclass[11pt]{article}
\usepackage[a4paper,margin=2.6cm]{geometry}
\usepackage[T1]{fontenc}
\usepackage[utf8]{inputenc}
\usepackage{newtxtext,newtxmath}
\usepackage{xcolor}
\usepackage[colorlinks=true,linkcolor=blue!60!black,urlcolor=blue!60!black]{hyperref}
\usepackage{enumitem}

\definecolor{quotegray}{gray}{0.35}
\newenvironment{reviewerquote}%
  {\medskip\noindent\color{quotegray}\itshape\begin{quotation}\noindent}%
  {\end{quotation}\medskip}
\newcommand{\resp}{\noindent\textbf{Response.}\ }
\newcommand{\changes}{\par\smallskip\noindent\textbf{Changes in the manuscript.}\ }
\newcommand{\secref}[1]{Section~#1}

\setlength{\parskip}{0.35em}

\begin{document}

\begin{flushleft}
{\large\bfseries Response to the reviewers}\\[0.4em]
Manuscript: \emph{Structure-Aware Row Scaling for Block-Structured Mixed-Integer Programs}\\
(submitted under the title \emph{Generator-Aware Positive Row Scaling for Block-Structured
Engineering Mixed-Integer Programs})\\
Author: Mathias Rodr\'iguez Castro\\
Manuscript ID: \texttt{[TO FILL]}
\end{flushleft}

\noindent Dear Editor and Reviewers,

Thank you for the careful and constructive evaluation. We have revised the manuscript
following every required (R) and suggested (S) point. The revision adds three new
experimental controls run on the same cluster and instance pools as the original study
(a solver-internal-scaling baseline, a hyperparameter sensitivity sweep, and an
operational stage ablation, now covering all three instance classes), plus a fourth
control that the ablation itself motivated (a role-blind ``flat'' scaling control), a
full original-scale equivalence audit of the four operational campaigns, and tempered
claims wherever the new evidence demanded it. All new code paths are additive and opt-in:
without the new flags the binaries are byte-identical to the version that produced the
published results (golden-file regression 23/23), so the original tables remain valid.
A marked manuscript (\texttt{latexdiff}) against the submitted version accompanies this
letter. The manuscript was retitled to \emph{Structure-Aware Row Scaling for
Block-Structured Mixed-Integer Programs}; the reproducibility repository was renamed
accordingly (\url{https://github.com/MathiasRodriguezCastro/structure-aware-row-scaling},
the old URL redirects) and its citation metadata updated.

\bigskip
\hrule
\section*{Required revisions}

\subsection*{R1 --- Solver-internal-scaling baseline}
\begin{reviewerquote}
Is the reported gain just a generic scaling effect that the solver could have produced
itself with its own scaling controls? [replace with verbatim comment]
\end{reviewerquote}
\resp We agree this is the decisive control and have run it in full: the unscaled model
under every admissible internal-scaling setting of each solver (Gurobi
\texttt{ScaleFlag} $0$--$3$; CPLEX \texttt{Read::Scale} $-1/0/1$), on the same instances,
for all three classes and both tolerances, compared against the structure-aware variants
under default internal scaling. Three findings. (i) Internal scaling is not a straw man:
disabling it is catastrophic (e.g.\ Simple Gurobi $\mathrm{sgm}$ $3.9\to42$\,s). (ii) On
the classes where the external rule is effective, no tuned internal setting reproduces
its runtime: the best internal Gurobi flag returns to the default level while the
structure-aware rule is $1.8\times$ faster (Simple), and no internal mode comes close on
SG-Ter-Mer. (iii) We also report honestly where this control does \emph{not} favor the
method: on the Full class at the strictest tolerance the retuned internal modes overtake
the external rule, and on Full under CPLEX the ``no scaling'' setting appears
dramatically faster but produces claimed-optimal objectives that disagree with the
default pipeline beyond the stopping tolerance on a third of the instances---we flag it
as numerically unreliable, which itself illustrates the paper's point that scaling
choices can invalidate certificates.
\changes New design paragraph in \secref{5.4} (additional controls); results in
\secref{6.7} with the extended Table~17 (three classes at $1\%$; the $0.1\%$ sweep
summarized in text); Discussion updated to scope the claim accordingly. Canonical CSVs
under \texttt{results-revision/r1-internal-scaling/} in the repository.

\subsection*{R2 --- Conditioning vs.\ runtime inconsistency}
\begin{reviewerquote}
SA-Aug is better conditioned but SA-Mat is often faster; conditioning does not line up
with runtime. [replace with verbatim comment]
\end{reviewerquote}
\resp We now state the variant-level inversion explicitly and re-frame conditioning as at
most a \emph{partial mediator} of the runtime effect throughout (abstract, results
lead-in, Discussion). The revision strengthens this further with two additions: the
correlation table now reports the per-variant decomposition---the pooled negative
conditioning--speedup association is carried by the Ruiz baselines
($\rho_S\approx-0.33/-0.36$), whereas for the structure-aware variants it is essentially
null---and the stage ablation (R5) shows that the stage that most improves the coupling
diagnostics is not the stage that drives the runtime. Precedent for both observations is
now cited (Elble \& Sahinidis 2012; Miltenberger, Ralphs \& Steffy 2018).
\changes \secref{5.6} (statistical analysis), Table~9 notes and surrounding text;
Discussion paragraph ``A mechanistic caveat\ldots''; new references.

\subsection*{R3 --- Hyperparameter sensitivity}
\begin{reviewerquote}
The constants $u=2$, $K=4$ and the coefficient window were fixed without justification.
[replace with verbatim comment]
\end{reviewerquote}
\resp The constants are now exposed as flags and swept on the Simple class under Gurobi:
$u\in\{1.5,2,4\}$, window $\in\{[10^{-6},10^{6}],[10^{-9},10^{9}],[10^{-12},10^{12}]\}$,
$K\in\{2,4,8\}$. The structure-aware rule is insensitive ($\mathrm{sgm}\in[2.01,2.12]$\,s
across the whole grid against a base of $3.73$\,s). We note transparently that $K$ is a
parameter of the Ruiz baselines only; varying it leaves the structure-aware rule
unchanged, as expected.
\changes \secref{5.4} (design) and \secref{6.9} with Table~19; canonical CSVs under
\texttt{results-revision/r3-hyperparam/}.

\subsection*{R4 --- Generality and the row-only premise}
\begin{reviewerquote}
The generality claims should be rescoped; the row-only choice is a premise, not a
finding. [replace with verbatim comment]
\end{reviewerquote}
\resp The abstract now states that all runtime evidence is single-platform and the method
is demonstrated on one hydrothermal testbed; a third limitation makes the row-only
premise explicit and identifies structured \emph{column} scaling as an untested companion
lever; the generalization to other domains is qualified by the shared structural premise
(sparse global coupling over local blocks). The title was kept accurate to the method,
with the scope made explicit in the text.
\changes Abstract; \secref{8} (limitations), third limitation; Conclusion.

\subsection*{R5 --- Operational stage ablation}
\begin{reviewerquote}
Local-only and coupling-only variants were never run on the real hydrothermal classes.
[replace with verbatim comment]
\end{reviewerquote}
\resp Both single-stage ablations now run on the real classes---including the Full
class---at both tolerances under Gurobi. The attribution is unambiguous and consistent:
the local stage alone recovers essentially the entire runtime benefit (Simple
$3.73\to2.01$\,s vs.\ $2.02$ for the full rule; Full $82.6\to70.9$ vs.\ $71.4$;
SG-Ter-Mer $14.4\to8.6$ vs.\ $8.6$), while the coupling stage alone is inert. We report
this even though it qualifies our own architecture, and we drew the natural consequence:
since the local kernel is computable from the flat matrix, we added a further
\emph{role-blind flat control} (the same per-row kernel and safeguards applied to every
row with no metadata at all) to isolate whether the row-role metadata is the active
ingredient of the runtime effect. \texttt{[TO FILL when the campaign completes: flat
control outcome and the corresponding scoping of the performance claim.]}
\changes \secref{5.4} (design), \secref{6.8} with Table~18 (three classes); Conclusion
updated (the previously listed ``future work'' ablation is now in the paper); flat
control code and campaign in the repository (\texttt{--plano}).

\subsection*{R6 --- Related-work positioning}
\begin{reviewerquote}
The related work should engage the classical equilibration literature and state precisely
what is new. [replace with verbatim comment]
\end{reviewerquote}
\resp The related work now engages geometric-/arithmetic-mean scaling (Tomlin),
Curtis--Reid, Sinkhorn--Knopp, Ruiz, and solver presolve (Achterberg et al.), and states
precisely what is new: not the kernel (the local factor is a geometric mean of per-row
extremes) but the metadata-driven \emph{source} of the scales and the staged
local--block--coupling \emph{ordering}. The revision adds Elble \& Sahinidis (2012) as
the empirical precedent that generic scaling payoffs are irregular, Miltenberger et al.\
(2018) on conditioning along the branch-and-cut tree, and Bergner et al.\ (2015) for
automatic structure detection as the fallback when the generator is not owned; the
absolute claim about modeling languages exporting flat matrices is now qualified
(block-structured extensions exist, but they feed decomposition algorithms rather than
numerical preprocessing---exactly the gap addressed).
\changes \secref{2.3}--\secref{2.4}; \secref{7} (deployability); references.

\section*{Suggested revisions}

\subsection*{S1 --- Deployability}
\resp New Discussion subsection \emph{Deployability and downstream effects}: metadata
availability in Pyomo/JuMP/GAMS/AMPL, instrumentation cost for owned generators, the
detection fallback (now cited), and a deliberately narrow value proposition.

\subsection*{S2 --- Dual variables}
\resp The same subsection now states that positive row scaling rescales constraint duals
by the inverse row factors, that the factors are stored, and that operators must descale
shadow prices (marginal costs, water values) before use.

\subsection*{S3 --- Proposition--proxy bridge}
\resp The propositions are explicitly downgraded to motivation, not guarantees; the
synthetic suite is identified as the empirical bridge, and the revision adds the
per-block audit trail via the controlled synthetic mechanism check.

\subsection*{S4 --- Abstract effect-size language}
\resp The abstract describes the Gurobi benefit as concentrated timeout reduction rather
than a broad per-instance speedup; the single-instance PAR10 caveat of SG-Ter-Mer is now
stated in the notes of \emph{both} tolerance tables.

\subsection*{S5 --- Partial correlations and missingness}
\resp The statistical section reports the size-controlled partial Spearman
($\rho\approx-0.25$, $n=1185$, essentially equal to the raw value), the balanced pools,
the $1$--$3\%$ per-run error rate, and the matched-pairs policy; the revision adds the
per-variant decomposition of the pooled correlation (see R2).

\section*{Editor's comments}

\subsection*{Length (W4)}
\begin{reviewerquote}
The manuscript should be shortened by roughly 15--20\%. [replace with verbatim comment]
\end{reviewerquote}
\resp \texttt{[TO FILL after the P2 editing pass: current status is that the required
controls added substantiating material; the dedicated trim pass consolidates the
local--block--coupling framing and the three-claims statements to single occurrences,
merges the two claim-summary tables, compresses the post-table narration paragraphs, and
moves the redundant 0.1\% Gurobi figures to Supporting Information.]}

\section*{New material added on our own initiative}
\begin{itemize}[leftmargin=1.4em]
  \item \emph{Original-scale equivalence audit} (\secref{6.1}, Table~6): the descaling
  verification promised by the design is now reported for all four operational campaigns,
  with verified explanations of every exceedance; notably, the audit shows the strict
  checker fails the \emph{unscaled base} at the same rate as the scaled variants where
  exceedances occur, and it catches wrong-optimal certificates in both directions.
  \item CPLEX tree-level \texttt{KappaStats}, so the numerical claim is now tabulated
  under both solvers.
  \item A run-to-run performance-variability paragraph in the limitations, using the
  repeated base campaigns as an empirical noise bound.
\end{itemize}

\bigskip
\noindent We believe the revision addresses every point and, where the new controls
qualified our own claims, we have reported those qualifications explicitly. We thank the
panel again for a review that materially improved the paper.

\bigskip
\noindent Sincerely,\\[0.3em]
Mathias Rodr\'iguez Castro\\
Facultad de Ingenier\'ia, Universidad de la Rep\'ublica, Uruguay

\end{document}
